% =========================================================================
% capitulos/01_aspectos_generales.tex
% Capítulo I: Introducción y Aspectos Generales de la Entidad
% =========================================================================

\section*{INTRODUCCIÓN}
\addcontentsline{toc}{section}{INTRODUCCIÓN}

La introducción debe presentar de manera clara y concisa el propósito del informe de prácticas preprofesionales, contextualizando el entorno en el que se llevaron a cabo las actividades. Se recomienda exponer brevemente la necesidad técnica o empresarial identificada en la institución receptora, el enfoque de ingeniería o metodología adoptado para su atención y los resultados principales alcanzados durante el periodo asignado.

Asimismo, es conveniente sintetizar la organización del documento indicando el contenido medular de cada uno de los capítulos: el Capítulo I comprende la descripción y aspectos generales de la entidad receptora; el Capítulo II expone el marco teórico y tecnológico de soporte; el Capítulo III detalla las actividades técnicas realizadas, objetivos, metodología y desarrollo; y finalmente, se presentan las conclusiones y recomendaciones derivadas de la experiencia profesional.

\section{Aspectos Generales de la Entidad}

\subsection{Información general de la institución}

\subsubsection{Identificación de la institución}
Describa la razón social, tipo de personería jurídica (pública o privada), número de Registro Único de Contribuyentes (RUC) y el rubro general de operaciones de la empresa o institución donde se ejecutaron las prácticas preprofesionales.

\subsubsection{Ubicación de la institución}
Indique la dirección física exacta de la sede u oficina principal (distrito, provincia y departamento). De ser pertinente, mencione las sucursales, portales web corporativos y canales oficiales de comunicación institucional.

\subsubsection{Misión y Visión}
\textbf{Misión:} \\
Describa la misión institucional establecida por la entidad receptora, enfocándose en su razón de ser, público objetivo y propuesta de valor en el mercado o sector en el que opera.

\vspace{0.3cm}
\textbf{Visión:} \\
Presente la visión estratégica de la entidad, orientada a sus metas de crecimiento, liderazgo y posicionamiento futuro en el mediano y largo plazo.

\subsubsection{Objetivos institucionales}
Enumere los fines primordiales de la entidad respecto a la prestación de sus servicios, comercialización de productos o gestión operativa:
\begin{itemize}
  \item Brindar soluciones y servicios especializados de alto estándar de calidad a sus clientes o usuarios finales.
  \item Fomentar la modernización y la transformación digital de sus procesos internos a través de herramientas tecnológicas eficientes.
  \item Promover la mejora continua, la sostenibilidad y la competitividad operativa dentro de su sector productivo.
\end{itemize}

\subsection{Actividades principales de la institución}
Detalle las líneas de negocio, servicios ofrecidos o funciones esenciales desempeñadas por la institución. Explique cómo estas actividades contribuyen al cumplimiento de su misión y cómo se articulan con las tecnologías de información.

\subsection{Área donde se realizaron las prácticas preprofesionales}
Especifique la gerencia, departamento u oficina técnica a la cual fue asignado el practicante (por ejemplo, Área de Desarrollo de Software, Unidad de Tecnologías de Información, etc.). Describa las responsabilidades del área, el cargo y grado del supervisor directo en la empresa y las dinámicas habituales de coordinación y control técnico.
